% A handout in the style of Edward Tufte, using the tufte-handout class.
% Sidenotes and margin figures sit beside the text they support.
\documentclass{tufte-handout}

\usepackage{amsmath}
\usepackage{booktabs}
\usepackage{microtype}

\title{{{title}}}
\author{{{author}}}
\date{\today}

\begin{document}

\maketitle

\newthought{Good graphics} tell the truth about data. They show the most
information with the least ink, and they leave the reader room to
think~\cite{tufte1983}.\sidenote{A sidenote is numbered and set in the margin
beside the line where it is called.}

\section{Data ink}

Every mark on a page should earn its place. The share of ink that carries
data is a simple measure of how much of a graphic does real work:
\begin{equation}
  r = \frac{\text{data ink}}{\text{total ink}}
\end{equation}
A high ratio means the reader sees the data first. Redundant gridlines,
heavy frames and decorative fills push the ratio down.
\begin{margintable}
  \centering
  \small
  \begin{tabular}{lr}
    \toprule
    Measure & Value \\
    \midrule
    Data ink & High \\
    Redundant ink & Low \\
    \bottomrule
  \end{tabular}
  \caption{A summary table set in the margin. Replace these rows with your own figures.}
\end{margintable}

\section{Tables}

Tables work best when the numbers are easy to compare down a column. Rules
are kept to the top, below the header, and at the bottom, and the rest of
the table stays open.\marginnote{This note has no number. Use one for an
aside that does not need a marker.}

\begin{table}[h]
  \centering
  \caption{Three books on information design, in order of publication.}
  \small
  \begin{tabular}{p{5.2cm}rp{3.4cm}}
    \toprule
    Book & Year & Theme \\
    \midrule
    The Visual Display of Quantitative Information & 1983 & Graphical excellence \\
    Envisioning Information & 1990 & Visual reasoning \\
    Beautiful Evidence & 2006 & Evidence and presentation \\
    \bottomrule
  \end{tabular}
\end{table}

\begin{fullwidth}
  Wide material, such as a long table or a sequence of panels, can run across
  the text column and into the margin. Use the full width when the content
  needs room to be read in one glance.
\end{fullwidth}

\section{Small multiples}

Repeating one small graphic with a shared scale lets the reader compare
groups at a glance~\cite{tufte2006}. A margin figure keeps such a graphic close to the text
that discusses it, without breaking the reading line.

\begin{marginfigure}
  \fbox{\parbox[c][3.2cm][c]{\dimexpr\linewidth-2\fboxsep-2\fboxrule}{\centering Placeholder for a margin figure}}
  \caption{A margin figure keeps its caption beside the text it supports.}
\end{marginfigure}

\bibliographystyle{plainnat}
\bibliography{references}

\end{document}
